%% ============================================================
%%  CHAPITRE 2 — Analyse des Besoins (Requirement Analysis)
%% ============================================================

\section*{Introduction}

Ce chapitre présente l'analyse des besoins du projet. À partir de la problématique identifiée
au chapitre précédent, nous identifions les acteurs impliqués, formulons les besoins fonctionnels
et non fonctionnels, modélisons les cas d'utilisation, et décrivons la méthodologie Agile
adoptée pour organiser la réalisation du projet en sprints itératifs.


\section{Identification des Acteurs}

Trois catégories d'acteurs interagissent avec la plateforme DevSecOps~:

\begin{table}[h]
    \centering
    \setlength{\arrayrulewidth}{0.8pt}
    \setlength{\dashlinedash}{3pt}
    \setlength{\dashlinegap}{2pt}
    \renewcommand{\arraystretch}{1.4}
    \caption{Acteurs du système et leurs rôles}
    \label{tab:acteurs}
    \begin{tabularx}{\textwidth}{|>{\centering\arraybackslash}X|>{\centering\arraybackslash}X|>{\centering\arraybackslash}X|}
\hline
\tableheader \textbf{Acteur} & \textbf{Type} & \textbf{Rôle} \\
\hline
Développeur & Humain & Écrit le code, pousse les commits, consulte les rapports de sécurité \\ \hdashline
Ingénieur DevOps / DevSecOps & Humain & Configure l'infrastructure, le pipeline et les politiques de sécurité \\ \hdashline
Pipeline CI/CD (GitHub Actions) & Système & Exécute automatiquement les builds, scans et déploiements \\ \hdashline
Argo CD & Système & Réconcilie l'état du cluster avec l'état désiré dans Git \\ \hdashline
Utilisateur final & Humain & Utilise l'application Todo App via le navigateur \\
\hline
    \end{tabularx}
\end{table}


\section{Besoins Fonctionnels}

Les besoins fonctionnels couvrent les capacités que la plateforme doit offrir. Ils sont
organisés par domaine dans le tableau~\ref{tab:besoins_fonctionnels}~:

\begin{table}[H]
    \centering
    \setlength{\arrayrulewidth}{0.8pt}
    \setlength{\dashlinedash}{3pt}
    \setlength{\dashlinegap}{2pt}
    \renewcommand{\arraystretch}{1.4}
    \caption{Besoins fonctionnels du projet}
    \label{tab:besoins_fonctionnels}
    \begin{tabularx}{\textwidth}{|>{\centering\arraybackslash}X|>{\centering\arraybackslash}X|>{\centering\arraybackslash}X|}
\hline
\tableheader \textbf{ID} & \textbf{Besoin Fonctionnel} & \textbf{Domaine} \\
\hline
BF-01 & Provisionner l'infrastructure GCP de manière reproductible via du code & Infrastructure \\ \hdashline
BF-02 & Construire et publier automatiquement les images Docker à chaque push & CI/CD \\ \hdashline
BF-03 & Exécuter des tests unitaires automatiquement à chaque commit & CI/CD \\ \hdashline
BF-04 & Scanner le code source pour détecter les vulnérabilités (SAST) & Sécurité \\ \hdashline
BF-05 & Scanner les dépendances npm et les images Docker pour les CVEs (SCA) & Sécurité \\ \hdashline
BF-06 & Scanner l'application déployée en boîte noire (DAST) & Sécurité \\ \hdashline
BF-07 & Déployer automatiquement l'application via GitOps (Argo CD) & Déploiement \\ \hdashline
BF-08 & Exécuter des tests E2E dans le cluster après chaque déploiement & Tests \\ \hdashline
BF-09 & Collecter et visualiser les métriques de l'application en temps réel & Observabilité \\ \hdashline
BF-10 & Surveiller les événements de sécurité en production (SIEM) & Sécurité \\ \hdashline
BF-11 & Permettre l'authentification, la gestion de tâches (CRUD) et la persistance & Application \\
\hline
    \end{tabularx}
\end{table}


\section{Besoins Non Fonctionnels}

Les besoins non fonctionnels définissent les contraintes de qualité de la plateforme~:

\begin{table}[H]
    \centering
    \setlength{\arrayrulewidth}{0.8pt}
    \setlength{\dashlinedash}{3pt}
    \setlength{\dashlinegap}{2pt}
    \renewcommand{\arraystretch}{1.4}
    \caption{Besoins non fonctionnels du projet}
    \label{tab:besoins_non_fonctionnels}
    \begin{tabularx}{\textwidth}{|>{\centering\arraybackslash}X|>{\centering\arraybackslash}X|>{\centering\arraybackslash}X|}
\hline
\tableheader \textbf{ID} & \textbf{Besoin Non Fonctionnel} & \textbf{Catégorie} \\
\hline
BNF-01 & Zéro clé d'authentification statique dans le pipeline (Keyless) & Sécurité \\ \hdashline
BNF-02 & Réseau Zero Trust~: tout trafic interdit par défaut, flux explicites & Sécurité \\ \hdashline
BNF-03 & Auto-scaling horizontal de 2 à 8 replicas selon la charge CPU & Performance \\ \hdashline
BNF-04 & Infrastructure 100\% reproductible à partir du code source & Fiabilité \\ \hdashline
BNF-05 & Chaque build, scan et déploiement doit produire un rapport auditable & Auditabilité \\ \hdashline
BNF-06 & Détection automatique et correction du drift de configuration (self-heal) & Résilience \\ \hdashline
BNF-07 & Couverture de tests unitaires $\geq$ 70\% & Qualité \\ \hdashline
BNF-08 & Temps de détection d'une CVE critique < 5 minutes & Sécurité \\ \hdashline
BNF-09 & Rollback automatique via Git revert + Argo CD & Disponibilité \\
\hline
    \end{tabularx}
\end{table}


\section{Diagramme de Cas d'Utilisation Global}

Le diagramme de cas d'utilisation global modélise les interactions entre les acteurs
et les principales fonctionnalités de la plateforme DevSecOps.

\begin{figure}[H]
    \centering
    \fbox{\parbox{0.85\textwidth}{\centering
        \vspace{3cm}
        Espace réservé pour la Figure~\ref{fig:use_case_global}\\
        \textbf{Diagramme de cas d'utilisation global}\\[6pt]
        \footnotesize Acteurs~: Développeur, Ingénieur DevSecOps, Utilisateur final,\\
        Pipeline CI/CD (système), Argo CD (système)\\[6pt]
        Cas d'utilisation~: Pousser du code, Construire \& tester, Scanner la sécurité,\\
        Déployer via GitOps, Surveiller, Gérer les tâches Todo
        \vspace{3cm}
    }}
    \caption{Diagramme de cas d'utilisation global de la plateforme DevSecOps}
    \label{fig:use_case_global}
\end{figure}


\section{Cas d'Utilisation Détaillés}

\subsection{CU-01~: Pousser du Code et Déclencher le Pipeline}

\begin{table}[H]
    \centering
    \setlength{\arrayrulewidth}{0.8pt}
    \renewcommand{\arraystretch}{1.4}
    \caption{Description du cas d'utilisation CU-01}
    \label{tab:cu01}
    \begin{tabularx}{\textwidth}{|>{\raggedright\arraybackslash}p{3.5cm}|X|}
\hline
\tableheader \textbf{Champ} & \textbf{Description} \\
\hline
Titre & Pousser du code et déclencher le pipeline \\
\hline
Acteur principal & Développeur \\
\hline
Précondition & Le développeur a un accès au dépôt Git \\
\hline
Scénario nominal & 1. Le développeur pousse un commit sur la branche principale \newline
2. GitHub Actions détecte le push et démarre le pipeline \newline
3. Le job \texttt{build-test} exécute lint, tests unitaires, SonarCloud et Trivy FS \newline
4. En cas de succès, le job \texttt{app-image-release} construit et publie l'image Docker \newline
5. Le job \texttt{update-k8s-tags} met à jour les manifests Kubernetes \\
\hline
Postcondition & L'image Docker est publiée, les manifests K8s sont mis à jour dans Git \\
\hline
Scénario alternatif & Si un test échoue ou un scan critique bloque, le pipeline s'arrête et notifie le développeur \\
\hline
    \end{tabularx}
\end{table}

\subsection{CU-02~: Déployer via GitOps}

\begin{table}[H]
    \centering
    \setlength{\arrayrulewidth}{0.8pt}
    \renewcommand{\arraystretch}{1.4}
    \caption{Description du cas d'utilisation CU-02}
    \label{tab:cu02}
    \begin{tabularx}{\textwidth}{|>{\raggedright\arraybackslash}p{3.5cm}|X|}
\hline
\tableheader \textbf{Champ} & \textbf{Description} \\
\hline
Titre & Déployer automatiquement via GitOps \\
\hline
Acteur principal & Argo CD (système) \\
\hline
Précondition & Les manifests K8s dans Git ont été mis à jour avec le nouveau tag d'image \\
\hline
Scénario nominal & 1. Argo CD détecte un changement dans le dépôt Git \newline
2. Argo CD compare l'état désiré (Git) avec l'état réel du cluster \newline
3. Argo CD synchronise automatiquement les ressources Kubernetes \newline
4. Les hooks PostSync lancent les tests Playwright E2E \newline
5. Le dashboard Argo CD affiche l'état synchronisé \\
\hline
Postcondition & L'application est déployée, les tests E2E sont validés \\
\hline
Scénario alternatif & En cas de drift détecté, Argo CD applique le self-heal automatiquement \\
\hline
    \end{tabularx}
\end{table}

\subsection{CU-03~: Gérer les Tâches Todo}

\begin{table}[H]
    \centering
    \setlength{\arrayrulewidth}{0.8pt}
    \renewcommand{\arraystretch}{1.4}
    \caption{Description du cas d'utilisation CU-03}
    \label{tab:cu03}
    \begin{tabularx}{\textwidth}{|>{\raggedright\arraybackslash}p{3.5cm}|X|}
\hline
\tableheader \textbf{Champ} & \textbf{Description} \\
\hline
Titre & Gérer les tâches Todo (CRUD) \\
\hline
Acteur principal & Utilisateur final \\
\hline
Précondition & L'utilisateur est authentifié sur l'application \\
\hline
Scénario nominal & 1. L'utilisateur ajoute une nouvelle tâche \newline
2. L'utilisateur coche une tâche comme complétée \newline
3. L'utilisateur supprime une tâche \newline
4. Les données persistent via localStorage après rechargement \\
\hline
Postcondition & Les tâches sont sauvegardées et visibles \\
\hline
Scénario alternatif & Si l'utilisateur n'est pas authentifié, il est redirigé vers la page de login \\
\hline
    \end{tabularx}
\end{table}


\section{Méthodologie de Travail~: Scrum Agile}

\subsection{Choix de la Méthodologie}

Le projet suit la méthodologie \textbf{Scrum} pour sa capacité à gérer la complexité
de manière itérative et incrémentale. Chaque sprint livre un incrément fonctionnel et
testable, permettant de valider progressivement les composantes de la plateforme DevSecOps.

\begin{figure}[H]
    \centering
    \fbox{\parbox{0.85\textwidth}{\centering
        \vspace{2.5cm}
        Espace réservé pour la Figure~\ref{fig:scrum_process}\\
        \textbf{Processus Scrum appliqué au projet}\\[6pt]
        \footnotesize Product Backlog $\to$ Sprint Planning $\to$ Sprint (2--3 semaines)\\
        $\to$ Sprint Review $\to$ Sprint Retrospective $\to$ Incrément livré
        \vspace{2.5cm}
    }}
    \caption{Processus Scrum appliqué au projet}
    \label{fig:scrum_process}
\end{figure}

\subsection{Équipe Scrum}

\begin{table}[H]
    \centering
    \setlength{\arrayrulewidth}{0.8pt}
    \setlength{\dashlinedash}{3pt}
    \setlength{\dashlinegap}{2pt}
    \renewcommand{\arraystretch}{1.4}
    \caption{Composition de l'équipe Scrum}
    \label{tab:equipe_scrum}
    \begin{tabularx}{\textwidth}{|>{\centering\arraybackslash}X|>{\centering\arraybackslash}X|>{\centering\arraybackslash}X|}
\hline
\tableheader \textbf{Rôle Scrum} & \textbf{Personne} & \textbf{Responsabilité} \\
\hline
Product Owner & Encadrant professionnel (Sopra Steria) & Définition des priorités et validation des livrables \\ \hdashline
Scrum Master & Encadrant académique (ESPRIT) & Suivi méthodologique et levée des blocages \\ \hdashline
Développeur & NAFFATI & Conception, implémentation et tests \\
\hline
    \end{tabularx}
\end{table}

\subsection{Planification des Sprints}

Le projet est découpé en \textbf{quatre sprints} de 2 à 3 semaines chacun, couvrant
l'ensemble des besoins identifiés~:

\begin{table}[H]
    \centering
    \setlength{\arrayrulewidth}{0.8pt}
    \setlength{\dashlinedash}{3pt}
    \setlength{\dashlinegap}{2pt}
    \renewcommand{\arraystretch}{1.4}
    \caption{Planification des sprints du projet}
    \label{tab:sprints}
    \begin{tabularx}{\textwidth}{|>{\centering\arraybackslash}X|>{\centering\arraybackslash}X|>{\centering\arraybackslash}X|>{\centering\arraybackslash}X|}
\hline
\tableheader \textbf{Sprint} & \textbf{Durée} & \textbf{Objectif} & \textbf{Besoins couverts} \\
\hline
Sprint 1 & 3 semaines & Infrastructure as Code et cluster GKE & BF-01, BNF-04 \\ \hdashline
Sprint 2 & 3 semaines & Containerisation, pipeline CI/CD et authentification Keyless & BF-02, BF-03, BNF-01, BNF-05 \\ \hdashline
Sprint 3 & 3 semaines & Sécurité intégrée (SAST, SCA, DAST), GitOps et réseau Zero Trust & BF-04, BF-05, BF-06, BF-07, BNF-02, BNF-06 \\ \hdashline
Sprint 4 & 2 semaines & Tests E2E, observabilité, SIEM et bilan & BF-08, BF-09, BF-10, BNF-03, BNF-07, BNF-08, BNF-09 \\
\hline
    \end{tabularx}
\end{table}

\subsection{Product Backlog}

Le Product Backlog regroupe l'ensemble des \textit{User Stories} priorisées par le
Product Owner. Le tableau~\ref{tab:backlog} présente les éléments principaux~:

\begin{table}[H]
    \centering
    \setlength{\arrayrulewidth}{0.8pt}
    \setlength{\dashlinedash}{3pt}
    \setlength{\dashlinegap}{2pt}
    \renewcommand{\arraystretch}{1.4}
    \caption{Product Backlog (extrait)}
    \label{tab:backlog}
    \begin{tabularx}{\textwidth}{|>{\centering\arraybackslash}X|>{\centering\arraybackslash}X|>{\centering\arraybackslash}X|>{\centering\arraybackslash}X|}
\hline
\tableheader \textbf{ID} & \textbf{User Story} & \textbf{Priorité} & \textbf{Sprint} \\
\hline
US-01 & En tant que DevOps, je veux provisionner l'infra via Terraform & Haute & 1 \\ \hdashline
US-02 & En tant que DevOps, je veux un cluster GKE Autopilot opérationnel & Haute & 1 \\ \hdashline
US-03 & En tant que Dev, je veux un Dockerfile multi-stage optimisé & Haute & 2 \\ \hdashline
US-04 & En tant que Dev, je veux un pipeline CI qui build et teste à chaque push & Haute & 2 \\ \hdashline
US-05 & En tant que DevSecOps, je veux une auth CI sans clé statique (WIF) & Haute & 2 \\ \hdashline
US-06 & En tant que DevSecOps, je veux des scans SAST/SCA automatiques & Haute & 3 \\ \hdashline
US-07 & En tant que DevSecOps, je veux un déploiement GitOps avec Argo CD & Haute & 3 \\ \hdashline
US-08 & En tant que DevSecOps, je veux un réseau Zero Trust (NetworkPolicies) & Moyenne & 3 \\ \hdashline
US-09 & En tant que DevSecOps, je veux un scan DAST post-déploiement (ZAP) & Moyenne & 3 \\ \hdashline
US-10 & En tant que QA, je veux 6 scénarios Playwright E2E in-cluster & Moyenne & 4 \\ \hdashline
US-11 & En tant que Ops, je veux Prometheus + Grafana pour l'observabilité & Moyenne & 4 \\ \hdashline
US-12 & En tant que SecOps, je veux Wazuh SIEM pour la surveillance runtime & Moyenne & 4 \\
\hline
    \end{tabularx}
\end{table}

\begin{figure}[H]
    \centering
    \fbox{\parbox{0.85\textwidth}{\centering
        \vspace{2.5cm}
        Espace réservé pour la Figure~\ref{fig:gantt}\\
        \textbf{Diagramme de Gantt du projet}\\[6pt]
        \footnotesize Sprint 1 (semaines 1--3) | Sprint 2 (semaines 4--6) |\\
        Sprint 3 (semaines 7--9) | Sprint 4 (semaines 10--11)
        \vspace{2.5cm}
    }}
    \caption{Diagramme de Gantt de la planification des sprints}
    \label{fig:gantt}
\end{figure}


\section*{Conclusion}

Ce chapitre a formalisé les besoins du projet à travers onze besoins fonctionnels et neuf
besoins non fonctionnels, identifié les acteurs du système, et modélisé les cas d'utilisation
clés. La méthodologie Scrum adoptée organise la réalisation en quatre sprints itératifs,
chacun livrant un incrément fonctionnel et testable. Le chapitre suivant présente l'analyse
et la conception détaillée du système.
